\printedpage{170}
% head: A. Grothendieck : Dérivateurs, Ch. XII, Édition des Universités de Montpellier II et Paris VII

All these results, and the notion of a $W$-cartesian square, dualize (using the appropriate axiom of right factorization, cf. page 240) into a theory of $W$-homotopic amalgamated sums, and of $W$-cocartesian squares. This time we work with the category
\[
\Phi \;=\; \left(
\begin{tikzcd}[ampersand replacement=\&, row sep=small, column sep=small]
 \& 1 \arrow[dl] \\
0 \& \\
 \& 2 \arrow[ul]
\end{tikzcd}
\right)
\qquad\quad \overset{j}{\hookrightarrow} \qquad\quad
\tilde{\Phi} \;=\; \left(
\begin{tikzcd}[ampersand replacement=\&, row sep=small, column sep=small]
 \& 1 \arrow[dl] \& \\
0 \& \& 3 \arrow[ul] \arrow[dl] \\
 \& 2 \arrow[ul] \&
\end{tikzcd}
\right),
\]
$j$ being a closed immersion. We have a restriction functor
\[
j^* : \Ho_W(\tilde{\Phi}) \longrightarrow \Ho_W(\Phi),
\]
for which we construct a left adjoint
\[
j_! : \Ho_W(\Phi) \longrightarrow \Ho_W(\tilde{\Phi}).
\]
And we find a natural notion of a
\msp{page 262}
$W$-cocartesian square, which seems irreproachable from every point of view, and of perfect simplicity. % [sic] tout points -> tous points, agreement of "tout"

This leads us to look again at the notion in question in Cat, where it looked very clever, deep and a little obscure, with integrals and all that. This is the moment to realize that ``that's all it is''. So to prove the

\textbf{Proposition 20.} \textit{Suppose} $\mathcal{M} = \mathrm{Cat}$\textit{, $W$ a basic localizer satisfying W(7 bis). Consider a commutative square in} Cat
\[
Q = \left(
\begin{tikzcd}[ampersand replacement=\&]
S \arrow[r] \arrow[d] \& X_1 \arrow[d] \\
X_2 \arrow[r] \& X_3
\end{tikzcd}
\right).
\]
\textit{For it to be $W$-cartesian} \texttt{[rather} $W$\texttt{-cocartesian]} \textit{in the sense of the theory developed here (for general categories of models), it is necessary and sufficient that}
\[
\varphi_Q : \quad \mathcal{X}_Q = \left(
\begin{tikzcd}[ampersand replacement=\&, row sep=small, column sep=small]
 \& X_1 \arrow[dl] \\
\int S \& \\
 \& X_2 \arrow[ul]
\end{tikzcd}
\right) \qquad\quad \longrightarrow \qquad\quad X_3
\]

\texttt{[rather}
\[
\varphi_Q : \quad \mathcal{X}_Q = \left(
\begin{tikzcd}[ampersand replacement=\&, row sep=small, column sep=small]
 \& X_1 \\
\int S \arrow[ur] \arrow[dr] \& \\
 \& X_2
\end{tikzcd}
\right) \qquad\quad \longrightarrow \qquad\quad X_3\text{\texttt{]}}
\]
\printedpage{171}
% head: A. Grothendieck : Dérivateurs, Ch. XII, Édition des Universités de Montpellier II et Paris VII
\textit{be in $W$.} (NB The interest of this criterion is that it is ``intrinsic'' to $Q$, not calling on a factorization of $\beta_1$ or of $\beta_2$ \texttt{[where} $\beta_i : S \longrightarrow X_i$\texttt{]} into $S \xrightarrow{\;\in \Cof_W\;} \tilde{X}_i \xrightarrow{\in W} X_i$. Cf. interesting application in section V 3.4., p. 146 ff., th. 4 and corollaries.)

\msp{page 263}
\textsc{Proof.} Consider a factorization of $S \longrightarrow X_1$ as
\[
S \xrightarrow{i_1} \tilde{X}_1 \xrightarrow{p_1} X, % [sic] X -> X_1, the factorisation is of S -> X_1
\]
with $i_1$ an open Dwyer immersion, $p_1 \in W$. So we have a commutative square
\[
\tilde{Q} = \left(
\begin{tikzcd}[ampersand replacement=\&]
S \arrow[r, hook, "\mathrm{Dw}"] \arrow[d] \& \tilde{X}_1 \arrow[d] \\
X_2 \arrow[r] \& X_3
\end{tikzcd}
\right)
\]
and a homomorphism $\tilde{Q} \xrightarrow{\alpha} Q$ which is the identity on the three vertices other than $\tilde{X}_1$, and which is $p_1$ at $\tilde{X}_1$. We deduce from this a commutative diagram
\[
\begin{tikzcd}[ampersand replacement=\&, row sep=small, column sep=small]
\mathcal{X}_Q \arrow[rr, "\mathcal{X}_\alpha = \lambda"] \arrow[dr, "\varphi_Q"'] \& \& \mathcal{X}_{\tilde{Q}} \arrow[dl, "\varphi_{\tilde{Q}}"] \\
 \& X_3, \&
\end{tikzcd}
\]
and we have $\lambda \in W$ (for it is a $\Phi$-arrow of fibered categories over $\Phi$, which induces a $W$-equivalence on each of three fibers.) Thus $\varphi_Q \in W \iff \varphi_{\tilde{Q}} \in W$, on the other hand it is clear that $Q$ is $W$-cocartesian if and only if $\tilde{Q}$ is. So we may assume that $S \hookrightarrow X_1$ is \texttt{[an]} open (Dwyer) immersion (\footnote[134]{Here we use W(3) (strong form), i.e. Loc(3 bis) in the notation of XVI. But we can do without it, cf. below.}). Applying this
\msp{page 264}
with the roles of $X_1$ and $X_2$ interchanged, we see that we may moreover assume that $S \longrightarrow X_2$ is also \texttt{[an]} open Dwyer immersion. Consider then
\[
\begin{tikzcd}[ampersand replacement=\&]
S \arrow[r, hook] \arrow[d, hook] \arrow[dr, phantom, "Q'"] \& X_1 \arrow[d, hook] \& \\
X_2 \arrow[r, hook] \& X_1 \amalg_S X_2 \arrow[r, "\pi"] \& X_3 \\
 \& \mathcal{X}_Q \arrow[u, "\varphi_{Q'}"] \arrow[ur, "\varphi_Q"'] \&
\end{tikzcd}
\]
\printedpage{172}
% head: A. Grothendieck : Dérivateurs, Ch. XII, Édition des Universités de Montpellier II et Paris VII
We have $\varphi_Q = \pi\varphi_{Q'}$, on the other hand we know that $\varphi_{Q'} \in W$ (cor. 3, p. 149). Thus
\[
\pi \in W \qquad\quad \iff \qquad\quad \varphi_Q \in W,
\]
where $\pi \in W$ means that $Q$ is $W$-cocartesian, q.e.d.

\msp{page 265}

\section*{8\quad Appendix : Return to the axioms W(7), W(7 bis) (pages 75, 76).}

(\footnote[135]{Provisional reflection, superseded by XVI.}). We have seen (prop. 6, p. 78) that W(7 bis) implies W(7). Let us prove a partial converse. % [sic] Reflexion -> Réflexion, missing accent; supplanté -> supplantée, agreement (in footnote 135)

\textbf{Proposition.} \textit{Suppose} W(7 bis) (cocartesian square axiom) \textit{satisfied, and let us place ourselves under the conditions of the statement of} W(7), \textit{and more generally, in the case of a cubical diagram}
\[
\begin{tikzcd}[ampersand replacement=\&, row sep=small, column sep=small]
 \& X_2 \arrow[rr] \arrow[dd, "f_2" near end] \& \& X_3 \arrow[dd, "f_3"] \\
X_0 \arrow[ur] \arrow[rr, crossing over] \arrow[dd, "f_0"'] \& \& X_1 \arrow[ur] \& \\
 \& Y_2 \arrow[rr] \& \& Y_3 \\
Y_0 \arrow[ur] \arrow[rr] \& \& Y_1 \arrow[ur] \arrow[from=uu, crossing over, "f_1"' near end] \&
\end{tikzcd}
\]
\textit{where the two squares $X$ and $Y$ (upper and lower face of the cube) are $W$-cocartesian. Consider the following conditions on this diagram}
\begin{enumerate}[label={}]
\item[\textit{a)}] \textit{$f_0 \in W$.}
\item[\textit{b)}] \textit{$f_1$ and $f_2 \in W$.}
\item[\textit{c)}] \textit{$f_3 \in W$.}
\end{enumerate}
\textit{Then we have the implications}
\[
a \text{ \textit{\&} } b \Longrightarrow c, \qquad\quad b \text{ \textit{\&} } c \Longrightarrow a.
\]
\textit{Moreover, if $W$ is stable under direct factors (e.g. $W$ strongly saturated), then we also have}
\[
a \text{ \textit{\&} } c \Longrightarrow b,
\]
\textit{thus the conjunction} \texttt{[of]} \textit{any two among the conditions a) b) c) implies the third.}

\textsc{Proof.} Using the functorial factorization
\printedpage{173}
% head: A. Grothendieck : Dérivateurs, Ch. XII, Édition des Universités de Montpellier II et Paris VII
\msp{page 266}
of an arbitrary arrow (here the $f_i$) as $pi$, with $i$ an open Dwyer immersion and $p \in W$ (p. 229), we are reduced to the case where the $f_i$ are open Dwyer immersions. But for any arrow $g : S \longrightarrow T$ which is in $\Cof_W$, let $Z(g)$ be its mapping cone % [sic] remené -> ramené, spelling
\[
\begin{tikzcd}[ampersand replacement=\&, row sep=small]
S \arrow[rr, "g"] \arrow[dr] \& \& T \arrow[dr, "\alpha(g)"] \& \\
 \& e \arrow[rr, "\bar{g}"] \& \& Z(g),
\end{tikzcd}
\]
and
\[
\alpha(g) : T = \mathrm{but}(g) \longrightarrow Z(g)
\]
the canonical arrow (\footnote[136]{This is a bit quick, for it is not clear that $X_1' \amalg_{X_0} X_2' \longrightarrow X_1 \amalg_{X_0} X_2$ is \texttt{[in]} $W$, without already using one of the three ingredients of W(7) (namely $f_0, f_1, f_2 \in W \implies f_3 \in W$). It seems then that our argument works only with an additional hypothesis of this sort \dots But there is \textit{Brown's lemma,} cf. XVI \S3, notably p. 20.}). This is functorial in $g$, so that the given cubical diagram extends to a double cube
\[
\begin{tikzcd}[ampersand replacement=\&, row sep=small, column sep=small]
 \& X_2 \arrow[rr] \arrow[dd, "f_2" near end] \& \& X_3 \arrow[dd, "f_3"] \\
X_0 \arrow[ur] \arrow[rr, crossing over] \arrow[dd, "f_0"'] \& \& X_1 \arrow[ur] \& \\
 \& Y_2 \arrow[rr] \arrow[dd, "\alpha_2" near end] \& \& Y_3 \arrow[dd, "\alpha_3"] \\
Y_0 \arrow[ur] \arrow[rr, crossing over] \arrow[dd, "\alpha_0"'] \& \& Y_1 \arrow[ur] \arrow[from=uu, crossing over, "f_1"' near end] \& \\
 \& Z_2 \arrow[rr] \& \& Z_3 \\
Z_0 \arrow[ur] \arrow[rr] \& \& Z_1 \arrow[ur] \arrow[from=uu, crossing over, "\alpha_1"' near end] \&
\end{tikzcd}
\]
with
\[
Z_i = Z(f_i),\ \alpha_i = \alpha(f_i) \qquad\quad \text{for } i \in \{0,1,2,3\}.
\]
Note that W(7 bis) implies that if $g \in \Cof_W$,
\msp{page 267}
we have
\[
g \in W \qquad\quad \iff \qquad\quad \bar{g} \in W, \qquad\quad \text{i.e. } Z(g)\ W\text{-aspheric},
\]
in particular
\[
f_i \in W \qquad\quad \iff \qquad\quad Z_i\ W\text{-aspheric}.
\]
Thus, the statement to be proved about the initial cube reduces to a statement about the square
\printedpage{174}
% head: A. Grothendieck : Dérivateurs, Ch. XII, Édition des Universités de Montpellier II et Paris VII
\[
\begin{tikzcd}[ampersand replacement=\&, row sep=small, column sep=small]
 \& Z_1 \arrow[dr] \& \\
Z_0 \arrow[ur] \arrow[dr] \& \& Z_3, \\
 \& Z_2 \arrow[ur] \&
\end{tikzcd}
\]
namely that the three conditions
\begin{enumerate}[label={}]
\item[a$'$)] $Z_0$ $W$-aspheric.
\item[b$'$)] $Z_1$ and $Z_2$ $W$-aspheric.
\item[c$'$)] $Z_3$ $W$-aspheric.
\end{enumerate}
giving rise to the implications
\[
\mathrm{a}' \text{ \& } \mathrm{b}' \Longrightarrow \mathrm{c}', \qquad\quad \mathrm{b}' \text{ \& } \mathrm{c}' \Longrightarrow \mathrm{a}',
\]
and that if $W$ is stable under direct factors, then
\[
\mathrm{a}' \text{ \& } \mathrm{c}' \Longrightarrow \mathrm{b}'.
\]

This is more or less the corollary on page 76, supplemented on page 77 (corollary), except that there the arrows are assumed to be open immersions, and the square cartesian (hence also cocartesian) in Cat, i.e. $Z_3 = Z_1 \cup Z_2$, $Z_0 = Z_1 \cap Z_2$. But
\msp{page 268}
what matters for the validity of this statement is (as one might suspect) that $Z_*$ be $W$-cocartesian. So we must first prove this.

\textbf{Lemma.} \textit{If the cubical diagram p. 265 is such that the faces $X$, $Y$ are $W$-cocartesian, and the $f_i$ ($i \in \{0,1,2,3\}$) $W$-cofibrations, then the square $Z$ of the cofibers of the $f_i$ is itself also $W$-cocartesian.}

This lemma follows from the theory of the ``unfolding of cubical diagrams'', relative to a category of models $(\mathcal{M}, W)$ admitting a theory of $W$-homotopic cofibers, which I really shall have to write! (But the existence for $(\mathrm{Cat}, W)$ of such a theory of $W$-homotopic cofibers, and hence of that of unfoldings of diagrams in $\mathcal{M}$ of type $(\boldsymbol{\Delta}^1)^I$, $I$ a finite index set, depends only on W(7 bis) - we do not have to call on W(7).)

\msp{page 269}
This ``unfolding'' is an infinite diagram all of whose infinite sequences, in the three directions of the coordinate axes, are ``suspension sequences''. I reproduce the useful part of the unfolding. % [sic] deploiement -> déploiement, missing accent